% page 330 du fac-simile (image p-329 du scan)
% bloc 160.0 x 242.0 mm ; marge gauche 8.0 mm
\pgc{-12.700mm}{0.000mm}{
\l{\cel{3}322.}
\l{}
\l{\sou{kuniklo}. (\sou{zool.}) Quadripedo, ek la klaso \textquotedbl{}vageri\textquotedbl{}, di qua}
\l{\cel{3}la gambi dopa esas kelke plu longa kam le avana, e qua ex-}
\l{\cel{3}kavas ter-truo quan lu habitas. - L.\cel{1}\sou{lepus cuniculus}.-ISL.}
\l{}
\l{\sou{kupo}. Vazo drinkala, expansita, kun pedo, ek bronzo, ar-}
\l{\cel{3}jento, oro, kristalo, e c. - EFIS.}
\l{}
\l{\sou{kupeo}. Veturo +kluza (klozita), di qua la korpo havas la formo}
\l{\cel{3}di berlino di qua la fako avana esabus tranchita. - DEFIRS.}
\l{}
\l{\sou{kupelo}. (\sou{tekn.}) Mikra vazo poroza, quan onu uzas por pur-}
\l{\cel{3}igar oro ed arjento per litargiro, qua fortiras la aloyuro}
\l{\cel{3}ed absorbesas kun lu da la parieti, dum ke la oro e la}
\l{\cel{3}arjento restas maso en la fundo di la vazo. - DEFIRS.}
\l{}
\l{\sou{kuperoso}. (\sou{patol.}) Morbo di la pelo di la vizajo, karakteriz-}
\l{\cel{3}ata da makuli redatra. - FIS.}
\l{}
\l{\sou{kuplar}. (\sou{tekn.})Juntar du peci de fero per charniri, rivet-}
\l{\cel{3}aguri, e c. - DEFI.}
\l{}
\l{\sou{kupleto}. En kansono plura-strofa, omna parto quan finas re-}
\l{\cel{3}freno. - EFIRS.}
\l{}
\l{\sou{kupolo}. (\sou{arkitekt}.)Mi-sfero, di qua la formo esas olta di}
\l{\cel{3}kupo subversita (kom ex.: la kupolo di Panthéon, en Paris;}
\l{\cel{3}la kupolo di la katedralo Santa Paulo, en London ; la}
\l{\cel{3}kupolo di la Parlamento di Usa, en Washington). -DEFIRS.}
\l{}
\l{\sou{kupono}. Parto de-tranchita, separita de ensemblo (ensemblo ek}
\l{\cel{3}rent-akto, karto di porcionizo nutrivala, tabakala, dum}
\l{\cel{3}indijo nacionala). - DEFIRS.}
\l{}
\l{\sou{kupro}. (\sou{kemio}). Korpo simpla, metalo redatra, plu gisebla}
\l{\cel{3}kam oro e kam arjento, tre duktila e martelagebla. -}
\l{\cel{3}\sou{Simbolo kemiala}\cel{1}: Cu.\cel{2}- DEFIS.}
\l{}
\l{\sou{kupulo}. (bot.)\cel{2}Asembluro ek braktei, quaze soldita ye sua}
\l{\cel{3}bazo, formacanta quaza kupo mikra, qua cirkondas la}
\l{\cel{3}floro e salias pinto cirkum la frukto. - DEFIS.}
\l{}
\l{\sou{kura}r. (\sou{netrans.}) I. (\sou{aludante homo, animalo)}. Irar, per sprin-}
\l{\cel{3}gi, plu rapide kam marche : I. esforcante devancar ulu ; 2.}
\l{\cel{3}esforcante atingar ulu ; 3. esforcante forirar de ulu. -FIS.}
\l{}
\l{\sou{kuracar}. (\sou{trans.}) (\sou{aludante mediko}). Agar a malado ye la ma-}
\l{\cel{3}niero quan lu judikas kom apta risanigor lu. - DEFIRS.}
\l{}
\l{\sou{kuracao}.\cel{2}Liquoro facita ek la shelo di oranji bitra. -}
\l{\cel{3}DEFIRS.}
\l{}
\l{\sou{kurajar}. (\sou{trans.}) I. Ne-angorar koram danjero. - II. Esar}
\l{\cel{3}energioza koram obstaklo vinkenda. - III. Esar psike sat}
\l{\cel{3}forta pri rezistar ulu, ulo. - DEFIRS.}
}
